\begin{table}[htbp]
\centering
\caption[Momentos estadísticos y correlación de PC1--PC3 del viento a 10 m - Subconjunto p90]
{Asimetría ($\gamma_1$) y curtosis ($\gamma_2$) de la PC1 del viento a 10 metros, y correlación de Spearman ($r$) entre PC1, PC2 y PC3 (con signo) y la vorticidad relativa a 850 hPa asociada a la fase (con signo, negativa en el Hemisferio Sur), para el subconjunto extremo (p90). El signo de la correlación depende de la convención de signo de cada EOF (Sección~\ref{subsec:analisis_pc}); lo relevante es que el lado del patrón dominante en que cae cada ciclón está asociado a su intensidad. Los valores significativos ($p<0.05$) se muestran en negrita.}
\label{tab:stat_p90_w10}
\begin{tabular*}{\textwidth}{@{\extracolsep{\fill}}llrrrrr}
\hline
\textbf{Región} & \textbf{Fase} & \textbf{Asimetría ($\gamma_1$)} & \textbf{Curtosis ($\gamma_2$)} & \textbf{PC1} & \textbf{PC2} & \textbf{PC3} \\
\hline
\textbf{SBR} & Incipiente       & 0.16  & -0.56 & \textbf{-0.53} & 0.08 & \textbf{-0.58} \\
             & Intensificación & -0.25 & -0.65 & \textbf{-0.45} & \textbf{-0.27} & \textbf{0.58} \\
             & Madurez         & 0.36  & -0.55 & \textbf{-0.44} & 0.19 & -0.12 \\
             & Decaimiento     & 1.15  & 3.95  & \textbf{-0.66} & \textbf{-0.37} & 0.20 \\
\hline
\textbf{LPB} & Incipiente       & 0.77  & 0.42  & \textbf{-0.43} & 0.08 & -0.17 \\
             & Intensificación & -0.88 & 0.96  & \textbf{-0.45} & \textbf{-0.38} & \textbf{-0.49} \\
             & Madurez         & -0.93 & 0.60  & \textbf{-0.37} & 0.04 & -0.23 \\
             & Decaimiento     & -0.20 & -0.18 & \textbf{-0.54} & \textbf{-0.37} & \textbf{-0.56} \\
\hline
\textbf{ARG} & Incipiente       & 0.29  & -0.14 & \textbf{-0.43} & \textbf{0.21} & -0.09 \\
             & Intensificación & -0.38 & -0.19 & \textbf{-0.32} & \textbf{-0.15} & \textbf{-0.38} \\
             & Madurez         & 0.41  & 0.25  & \textbf{-0.26} & \textbf{-0.32} & \textbf{-0.33} \\
             & Decaimiento     & 0.51  & 1.66  & \textbf{-0.74} & -0.04 & \textbf{-0.20} \\
\hline
\end{tabular*}
\end{table}
